\scanpage{5}
\uline{NB}\quad Si $\sigma$ est un $\underset{\hspace{-1.3em}\unsure{\cdot}}{\text{automorphisme}}$ du corps des complexes, et si $f$ est ``spécial'', alors $f^\sigma$ est \unsure{égalmt|également} spécial -- donc $\operatorname{Aut}(\mathbb C)$ opere sur l'ens des $f$ spèciaux. Mais il opère par l'intermédiaire de $\operatorname{Aut}(\unsure{\overline{\mathrm Q}|\overline{\mathbb Q}})$, car si $f=P/Q$ spécial, les $\underset{\hspace{2.5em}\unsure{-}}{\text{coeff.}}$ de $P,Q$ sont des \uline{nbs algébriques}. Il revient au $=$ de dire que les zéros de $P,Q$ sont des nbs algébriques, i.e. le diviseur de $f$ est
% Isolated point below the middle of automorphisme: visible shape retained, intention uncertain (full scan x854-858, y168-174).
% Faint short trace below the end of coeff.: visible shape retained, deliberate writing uncertain (full scan x410-434, y892-903).

\[
 \unsure{-|=}\quad\unsure{-}\hspace{5em}-
\]
% The second, very faint short trace is near full-scan (430,1423); its intention is uncertain.
% archive sheet number: 4
